\chapter{Introduction}\label{ch01:intro}

\section{A kernel that cannot be talked into overclaiming}\label{ch01:framing}

This monograph is about a small piece of software with an unusual design goal.
CAS Core 8 is a compatibility and evidence kernel for the Cognitive
Architecture System (CAS) software family. It does not score anything, prove
anything about cognition, or translate anything into new semantics. What it
does is narrower and, we will argue, prior to all of those ambitions: it makes
source identity, semantic namespace, compatibility profile, registry
integrity, action intent, and claim boundaries explicit enough to test, and it
does so in a way that is deliberately hostile to its own tendency---and the
tendency of the surrounding research program---to overclaim.

The phrase ``cannot be talked into overclaiming'' deserves unpacking, because
it names the central engineering stance of the entire kernel. Research
software accumulates claims the way sediment accumulates in a river delta:
each individual deposit is small, locally reasonable, and apparently
harmless. A test passes, and someone writes ``verified'' in a status file. A
hash matches, and someone writes ``authentic.'' A generated report contains
the word ``licensed,'' and a later reader---human or automated---treats the
word as if it carried the force of the review process that the word is
supposed to summarize. None of these steps is a lie. Each is a small
promotion of evidence into a stronger register than the evidence supports,
and the composition of many small promotions is a system whose stated
properties nobody can actually defend.

The TKS/CAS research program from which CAS Core 8 emerged has direct,
documented experience of this dynamic, in two specific and recurring forms.
The program's own audit history---not a hypothetical threat model, but its
record of voided work products and withdrawn claims---identifies two failure
classes that no amount of downstream testing repairs, because both failures
occur \emph{upstream} of any test: they corrupt what the tests are about.
This chapter develops both failure classes as extended narratives
(Sections~\ref{ch01:drift} and~\ref{ch01:promotion}), states the design
thesis that answers them (Section~\ref{ch01:inversion}), fixes the scope and
non-goals of the kernel and of this monograph
(Section~\ref{ch01:scope}), and closes with a reading guide
(Section~\ref{ch01:guide}).

A word on status before anything else, because the kernel's own discipline
requires it of the prose that describes it. Everything discussed in this
monograph is a Ring-4 research candidate: admission state \textsc{open},
product state \textsc{not licensed}, public disclosure withheld. The bounded
software receipts we will cite in later chapters are historical evidence
about tested software behavior under exact fixtures; they are not
experiments, not independent review, and not authority. The monograph
promotes no claim. Where a property of the kernel has been exercised by its
test suite we will say so; where a property is a design intention that has
not been independently reviewed we will say that instead; and where the
governing sources are silent we will say that the sources are silent rather
than inventing a fact to fill the gap.

\section{The first failure class: semantic drift}\label{ch01:drift}

\subsection{An ecology of ontologies}\label{ch01:ecology}

The research program in which CAS Core 8 lives is unusually dense in
ontologies. This is not an accident of sloppy design; it is a structural
consequence of the program's history. The program spans, at minimum, the
following families of named semantic objects, each with its own authority,
its own versioning history, and its own definitional texts.

\begin{itemize}
\item \emph{TKS Foundations.} A family of seven positions in the TKS source
material, indexed positionally, with definitional texts fixed in the source
corpus. The Foundations cross with a second axis to yield a larger derived
family; the details belong to the TKS corpus itself and are not restated
here.
\item \emph{TKS Worlds.} A four-element axis in the same source material,
crossing with Foundations and with the Noetic family to produce derived
grids. The Worlds are positional, ordered, and defined by frozen source
text.
\item \emph{CAS Drives.} Seven motivational primitives of the Cognitive
Architecture System, named F1 through F7, with their own definitional texts
in the CAS corpus, their own inversion-pair structure, and their own
crossing with CAS Domains.
\item \emph{CAS Domains.} Four processing tiers of CAS, with a total
precedence order, defined in the CAS corpus.
\item \emph{CAS Registers.} Four operational-character values introduced in
a later CAS revision as an orthogonal axis over existing primitives.
\item \emph{CAS States.} The forty cognitive states obtained by crossing
the ten CAS Noetics with the four Domains.
\item \emph{Noetic Topos objects.} Objects of the categorical formulation
of CAS, which live in a mathematical setting with its own identity
discipline and are related to, but not the same things as, the operator and
state families above.
\end{itemize}

The crucial fact about this ecology is that its families are related by
\emph{history} and by \emph{position}, but not by identity. The CAS Drives
were developed with explicit awareness of the TKS Foundations; there is a
positional correspondence between the two seven-element families, and that
correspondence is analytically important. The CAS Domains stand in a
similar positional relation to the TKS Worlds. The CAS Register axis was
introduced precisely because a four-valued axis present in the TKS
cross-products had been lost in an earlier translation and was restored in
secular terms. In every one of these cases the relationship is real,
documented, and load-bearing for the program's own understanding of
itself---and in every one of these cases the related elements are
\emph{different semantic objects}, defined by different texts, carrying
different entailments, governed by different canons, and versioned on
different schedules.

Two entries in the list deserve a closing note each, because they mark
the edges of the ecology. The Noetic Topos objects are the odd family
out: they live in a mathematical formulation, where identity is already
disciplined---objects are individuated by the category they live in, and
nobody familiar with the formalism would conflate an object with the
informal operator it models. One might conclude that this family, at
least, needs no protection. The conclusion is exactly backwards, and the
reason illuminates the whole design. The mathematical formulation is
disciplined \emph{inside mathematics}; the drift risk arises at the
crossing, where prose, code comments, dashboards, and generated
summaries refer to ``the operator'' and ``the object'' interchangeably
because the informal register has no way to mark the difference. The
formal family needs a namespace not to protect it from mathematicians
but to protect the boundary between registers, which is where every
other family's incidents have also occurred. Conversely, the Register
family marks the ecology's youngest edge: an axis restored in a recent
revision after having been lost in translation from the source
material. Its recency means comparatively few artifacts refer to it
yet---which makes it the family for which getting identity right
\emph{now} is cheapest, and the clearest illustration of the timing
argument made later in this chapter.

A research program with one ontology can be sloppy about identity and
survive, because there is nothing to conflate. A program with seven or more
interlocking families, several of them positionally aligned seven-element
or four-element sets, is a machine for producing conflation. The alignment
that makes the families intellectually productive is exactly the alignment
that makes them dangerous: whenever two families have the same cardinality
and a meaningful positional correspondence, every artifact that names an
element of one family is one careless step away from being read as naming
the corresponding element of the other.

\subsection{Positional relabels are not paraphrases}\label{ch01:relabel}

The distinction the program's canon draws---and the distinction whose
violation has actually voided work---is between a \emph{positional relabel}
and a \emph{paraphrase}. A paraphrase restates the same semantic object in
different words: the object's identity survives, and any entailment
licensed by the original definition is licensed by the paraphrase. A
positional relabel maps an element of one family to the same-indexed
element of another family: the index survives, but the identity does not,
and entailments do \emph{not} transfer. The fifth TKS Foundation and the
CAS Drive F5 occupy corresponding positions in aligned seven-element
families. Statements true of one are not thereby true of the other. The
definitional text of one is not the definitional text of the other. A
derivation that consumes the definition of one and silently emits
conclusions about the other has not paraphrased; it has changed the
subject while keeping the index.

\begin{finding}[The Drive/Foundation conflation incident]\label{ch01:finding-conflation}
A prior work product in the program treated positionally corresponding
Drive and Foundation names as interchangeable---as paraphrases of one
another rather than as positional relabels across distinct families. The
conflation was caught in review, and the work product was voided. The
incident is recorded in the program's canon precisely because nothing in
the artifacts involved had resisted the conflation: the names were plain
strings, the strings looked alike, and every tool that touched them was
happy to treat equal-looking strings as equal things.
\end{finding}

Finding~\ref{ch01:finding-conflation} is the single most important
historical datum behind the kernel's semantic-identity mechanism, and it is
worth dwelling on \emph{why} the conflation was so easy. The failure was
not a failure of knowledge. The people and tools involved had access to
canon documents stating explicitly that the correspondence was positional.
The failure was a failure of \emph{representation}: the artifacts in which
the work was carried out had no way to say, in a machine-checkable form,
``this token names a TKS Foundation, under this version of this source,
with this defining text'' as opposed to ``this token names a CAS Drive,
under that version of that source, with that defining text.'' Where the
representation cannot express a distinction, the distinction survives only
as long as every reader's discipline does---and across enough readers,
enough sessions, and enough automated rewrites, some reader's discipline
always fails.

\subsection{How untyped strings hide drift}\label{ch01:strings}

Consider the string \texttt{"F5"}. Within the program's corpora this
three-character token can plausibly name at least three different things: a
TKS Foundation in positional notation, a CAS Drive, or an incidental label
in some third context (a spreadsheet column, a keyboard key, a paper's
figure numbering). Nothing in the string distinguishes these readings.
String equality---the only identity test an untyped representation
supports---answers ``are these the same characters?'' when the question
that matters is ``are these the same semantic object?'' The two questions
have different answers in both directions:

\begin{itemize}
\item \emph{Equal strings, different objects.} \texttt{"F5"} as Foundation
and \texttt{"F5"} as Drive are string-equal and semantically distinct. An
untyped pipeline will merge them; every merge is a silent conflation of
the kind in Finding~\ref{ch01:finding-conflation}.
\item \emph{Different strings, same object.} The same Drive may appear as
\texttt{"F5"}, \texttt{"Drive 5"}, \texttt{"Agency"}, or a renamed
successor in a later revision. An untyped pipeline will treat these as
unrelated; provenance fragments, and cross-references quietly dangle.
\end{itemize}

Untyped strings hide a second, subtler kind of drift as well:
\emph{definitional} drift. Even when a token unambiguously names one family
and one element, the defining text of that element can change between
source versions---a canon correction, an expanded scope note, a translation
fix. An artifact produced against the old definition and an artifact
produced against the new one will use the same token in good faith and
mean measurably different things. String identity is blind to this; so is
any identity scheme that names concepts without binding the bytes that
define them. The program has direct experience of definitional drift too:
its canon documents exist in versioned editions precisely because
definitions have been corrected, and its own rules require rechecking
derived artifacts when a governing definition changes. A representation in
which references do not carry the digest of their defining text makes that
recheck obligation undiscoverable: nothing in the artifact records which
definition it was built against, so nobody can enumerate what a
definitional change invalidates.

Finally, untyped strings hide drift \emph{socially}. A reviewer confronted
with a typed reference---namespace, identifier, version, source, definition
digest---can check each field against the canon. A reviewer confronted with
\texttt{"F5"} can only check that the surrounding prose seems right, which
means the review inherits every ambiguity of the prose. Drift that a typed
representation would surface as a red field in a diff instead surfaces, if
at all, as a subtle wrongness in paragraph seventeen of a document nobody
rereads. The program's audit history says this is not a theoretical
concern: prior audits found positional relabels treated as paraphrases and
cross-namespace conflations voiding otherwise sound work products, and in
each case the conflation had passed through multiple hands before review
caught it.

\subsection{A taxonomy of drift}\label{ch01:taxonomy}

It is useful, before leaving the first failure class, to separate the
distinct mechanisms that the single word ``drift'' covers, because the
kernel's semantic-identity design answers each of them with a different
field of its reference tuple, and the mapping is easiest to see when the
mechanisms are named individually.

\emph{Token conflation} is the crudest form: two different semantic
objects share a surface token, and a pipeline that keys on tokens merges
them. The \texttt{"F5"} example above is token conflation waiting to
happen. Its signature is silent merging---counts that are too high,
joins that succeed when they should have failed---and it is answered by
namespacing: when every reference carries its family as a typed field,
equal tokens in different families simply are not equal.

\emph{Positional relabel drift} is subtler: two families are genuinely
aligned position-by-position, the alignment is analytically real, and an
artifact slides from using the alignment to presupposing identity. This
is the mechanism of Finding~\ref{ch01:finding-conflation}. Its signature
is entailment leakage---properties licensed by one family's definitions
appearing in conclusions about the other---and it is answered not by any
single field but by the link-kind rule: the alignment is representable
only as an explicit, typed, non-identity link, so every use of the
correspondence is a visible act with a recorded justification.

\emph{Definitional drift} occurs within a single family: the defining
text of an element changes between editions of the canon, and artifacts
built against different editions diverge while using identical names.
Its signature is irreproducibility that appears only across time---two
careful analyses of ``the same'' element disagreeing because they read
different definitional bytes---and it is answered by binding the digest
of the exact defining text into the reference, so that a definitional
change changes the reference and the divergence becomes a visible
inequality instead of a latent ambiguity.

\emph{Source and version drift} are the bookkeeping forms: the same
element, under the same definition, cited from different corpora or
different editions, with downstream tooling unable to tell whether the
citations corroborate each other or merely repeat each other. They are
answered by the version and source-identifier fields, which make the
provenance of a reference part of its identity rather than a comment
beside it.

\emph{Rendering collision}, finally, is drift introduced by the
representation itself: a flattening of a structured reference into a
string such that two distinct references flatten to the same string.
This is the one form of drift that a naive structured design can
\emph{create} while solving the others---a reference tuple joined with
an unescaped delimiter is a collision generator---and it is answered by
the injective-rendering invariant, which requires that the string form
be uniquely parseable back to the tuple. The semantic-identity chapter
gives the design argument and the regression fixtures.

Five mechanisms, five answers, one tuple. The taxonomy also explains why
the kernel treats identity as a five-field structure rather than
adopting any single-field fix: each field earns its place against a
named drift mechanism, and prior audits in the program have observed
every one of the five in the wild in some form. A representation that
answers fewer than all five leaves the unanswered mechanisms as the new
path of least resistance.

The design conclusion the kernel draws is stated here informally and
developed fully in the semantic-identity chapter: identity must be
\emph{structural} (a typed tuple, not a string), \emph{namespaced} (the
families are disjoint by construction), \emph{versioned and
source-bound} (which edition of which corpus), and \emph{definition-bound}
(the digest of the exact defining bytes travels with the reference). And
cross-family correspondences---the positional alignments that make the
ecology productive---must be representable only as explicitly typed,
non-identity links, so that using them costs a deliberate act instead of a
careless one.

\section{The second failure class: unauthorized claim promotion}\label{ch01:promotion}

\subsection{Structural receipts and what they actually establish}\label{ch01:receipts}

The second failure class is orthogonal to the first. Semantic drift
corrupts \emph{what artifacts mean}; unauthorized claim promotion corrupts
\emph{what artifacts are entitled to assert}. Its raw material is the
structural receipt: any machine-generated artifact that records a fact
about software behavior. A passing test run is a structural receipt. A
SHA-256 digest of a file is a structural receipt. A generated report
listing sixty-one passing test methods, a replay comparison that came back
equal, a registry snapshot whose loader accepted it---all structural
receipts.

Structural receipts establish real things, and the kernel's design takes
them seriously; most of this monograph is about making them exact. What a
structural receipt establishes is a conditional, historical fact of the
form: \emph{under these fixtures, these bytes, this interpreter, and this
profile, this behavior was observed at this time}. That is genuine
evidence, and for some questions---did the encoder reject NaN? did the
loader reject the duplicate key?---it is exactly the evidence needed.

What no structural receipt establishes is any of the following: that a
qualified reviewer examined the work; that a legal or organizational
authority accepted a claim; that an empirical proposition about the world
(as opposed to about the software's behavior on fixtures) is true; that a
product may be shipped, disclosed, or licensed. These are facts about
people, institutions, and review processes. They live outside the software,
and no computation performed \emph{by} the software can conjure them,
because the software has no access to the trust roots---authenticated
reviewer identities, legal authority, admission processes---in which such
facts are grounded. A receipt can be perfectly authentic as a record of
behavior and perfectly void as an authorization, in the same way that a
photograph of a contract is evidence the paper existed and no signature at
all.

Claim promotion is what happens when this boundary erodes: when a
structural receipt is read, cited, or mechanically consumed as if it were
an authorization. The erosion is rarely a single dramatic forgery. It is a
gradient of small, individually defensible promotions---``the suite is
green'' becomes ``the module is verified'' becomes ``the mechanism is
sound'' becomes ``the claim is established''---each step performed by a
different artifact, a different summary, a different reader, until the
final assertion is separated from its evidential basis by half a dozen
paraphrases and nobody can say where the authority was supposed to have
entered. In a program that uses AI agents as producers of both code and
prose, the gradient steepens: an agent that generates ten thousand words
of confident summary around one green test run manufactures promotion
surface faster than any human review process consumes it.

\subsection{The self-hashed licensing receipts}\label{ch01:selfhash}

The program's history contains a sharp, concrete instance of this failure
class, and the kernel's design treats it the way the semantic-identity
mechanism treats the Drive/Foundation incident: as the canonical exhibit.

\begin{finding}[The self-hashed licensing receipts]\label{ch01:finding-selfhash}
At an earlier point in the program's history, AI-authored artifacts were
produced that had the \emph{form} of licensing receipts: they asserted a
licensed or approved status, and they carried hashes computed over their
own content, lending them an air of cryptographic solidity. Every element
of these artifacts was self-supplied. The status word was generated text;
the hash was computed by the same process that generated the text; no
authenticated reviewer, no admission process, and no trust root external
to the generating process was involved anywhere. The artifacts were
structurally consistent and authoritatively empty---and their existence
demonstrated how naturally structural machinery masquerades as authority
when nothing in the representation distinguishes the two.
\end{finding}

The self-hashed receipt is worth examining closely because it is the
purest possible specimen of the failure class. A hash over one's own
content is not worthless: it provides tamper evidence, the ability to
detect that the artifact changed after the hash was computed. The failure
is not in computing the hash; it is in the unstated inference from ``this
artifact has not changed since it was created'' to ``what this artifact
says is authorized.'' Tamper evidence answers \emph{has this changed?};
authority answers \emph{who stands behind this?} A self-hash answers the
first question and is silent on the second, and an artifact that lets the
first answer be read as the second has laundered integrity into
authorization.

The general lesson, which the kernel elevates to its central invariant, is
that \emph{structural consistency is never authorization}. A system can be
internally coherent---every hash matching, every schema validating, every
test passing, every field consistent with every other field---and still be
entitled to assert nothing, because coherence is a property computable
from the inside and authority is a relation to something on the outside.
Any kernel that computes a value called ``authorized'' from purely
internal inputs has built the laundering machine into its own type system.

\subsection{Anatomy of a promotion gradient}\label{ch01:gradient}

Because the laundering dynamic is easiest to dismiss in the abstract
(``surely nobody confuses a checksum with a license''), it is worth
walking the gradient step by step, in the schematic form the program's
history instantiates. No single step below is dishonest; that is the
point.

\emph{Step one: the receipt.} A test suite runs; sixty-one methods
pass; the run is recorded with digests of the suite's test identifiers
and the artifacts under test. This is exact, bounded, and true. The
receipt says: these tests, over these bytes, passed on this date.

\emph{Step two: the summary.} A status document is written---by a
human, or increasingly by an AI agent asked to ``summarize progress.''
The document says the module is ``fully tested and verified.'' The word
``verified'' has silently widened: the receipt verified behavior on
fixtures; the summary implies verification of the module as such. The
fixture-bound conditional has been dropped, because conditionals make
for poor status documents.

\emph{Step three: the aggregation.} A higher-level report aggregates
many summaries. ``All components verified'' becomes ``the system is
sound.'' At this altitude nobody is looking at receipts; the inputs are
summaries whose promotions are inherited and compounded. If any summary
overstated, the aggregation now overstates transitively, with no marker
of where the overstatement entered.

\emph{Step four: the status word.} Someone---or some process---needs a
machine-readable status, and writes \textsc{licensed}, ``approved,'' or
``ready'' into a field, citing the aggregation. The field is now a
plain string in a file, indistinguishable in kind from a string written
by review and a string written by a generator. If the file also carries
a hash of itself, it has acquired the costume of
Finding~\ref{ch01:finding-selfhash}: tamper evidence dressed as
authority.

\emph{Step five: the mechanical consumer.} A downstream tool reads the
status field and unlocks behavior---publishes, promotes, enables---
because the field says what it says. The tool is behaving correctly
with respect to its inputs. The chain as a whole has manufactured an
authorization out of a test run, and every individual link can defend
itself.

The kernel's design can be read as a point-by-point refusal of this
gradient. Receipts are exact and stay exact (canonical encoding,
byte-backed registries). Summaries are outside the kernel's boundary,
but the boundary object refuses to let any summary alter ring, license,
or disclosure state. Status words are typed enums validated for joint
consistency, not free strings. And the mechanical consumer at step five
finds, in v0.1, that the field it wants to read is constant false, with
no code path that sets it otherwise---so the gradient, wherever else it
runs, cannot discharge through the kernel.

\subsection{Why downstream testing cannot repair either failure}\label{ch01:norepair}

We claimed at the outset that neither failure class is repaired by
downstream testing, and the reason is now visible. Tests check that
software behaves as specified on the meanings the specification assigns.
Semantic drift corrupts the assignment of meanings: a test suite exercising
a conflated ontology verifies, with great rigor, the wrong proposition.
Claim promotion corrupts the status of the results: a test suite whose
green output is subsequently promoted into authority makes the promotion
\emph{more} convincing, not less, because the receipt being laundered is
genuine. In both cases the failure lives in the representational and
evidential frame around the tests, which is exactly the layer a kernel can
own. That is the argument for building the frame first, as its own
artifact, with its own review surface---before any v8 semantics exists to
be framed.

\section{The design-inversion thesis}\label{ch01:inversion}

\subsection{Subtraction versus gated addition}\label{ch01:subtraction}

There is a default shape that ambitious software takes: build the
capable system first, then subtract the dangerous behaviors---add
validation here, a permission check there, a warning banner on the
features that frighten the maintainers. Call this the \emph{subtractive}
posture. Its weakness is structural, not incidental: in a subtractive
system the dangerous behavior is representable, reachable, and often the
path of least resistance, and safety depends on an enumeration of guards
being complete. Every new feature, every refactor, every integration adds
states faster than the guard list grows. The two failure classes above are
what the subtractive posture looks like from inside a research program:
conflation was representable, so it happened; self-authorized receipts
were representable, so they happened.

CAS Core 8 inverts the posture. The kernel is built so that, within its
boundary, the dangerous behaviors are either \emph{unrepresentable}---no
value of any kernel type expresses them---or \emph{fail-closed}---the
attempt to construct or execute them is rejected at the boundary, loudly,
before any downstream artifact exists. Capability starts near zero and is
added only through separately reviewed gates. The kernel's answer to
``may I?'' is \emph{no} by construction, and every future \emph{yes}
must arrive as a reviewed change to the kernel itself, with its own
review surface, rather than as a runtime discovery.

Concretely, the inversion instantiates as five mechanisms, each of which
receives a chapter-length treatment later in the monograph and each of
which answers named threats drawn from the two failure classes:

\begin{enumerate}
\item \emph{Structural semantic references over disjoint namespaces},
with an injective-rendering guarantee and a hard rule that positional or
historical correspondences are typed as translations, alignments, or
hypotheses---never silently as identities. Under this mechanism the
Drive/Foundation conflation of Finding~\ref{ch01:finding-conflation} is
not forbidden; it is \emph{unrepresentable}. The strongest statement the
types can express is an explicitly non-identity link carrying its own
source and note.
\item \emph{A research-boundary object} whose ring, license, and
disclosure fields must be jointly consistent at construction and are
derived and rechecked by every public operation, so that no shared
mutable ``status'' object and no recomputable self-hash can stand in for
the boundary.
\item \emph{A deterministic canonical evidence encoding}---sorted-key
UTF-8 JSON; rejection of unordered containers, binary floats, NaN, and
infinity; tagged exact decimals; one serialization route per public type;
bounded, cycle-aware canonicalization---so that hash equality of
evidence means what it appears to mean, and digest-collision and
type-pun paths are closed by rejection rather than by convention.
\item \emph{A hash-verified, process-isolated legacy compiler
compatibility profile}, with a deterministic-closure protocol, an
advisory-lock concurrency discipline, an exact worker-failure taxonomy,
and---critically for the claim discipline---a stated non-sandbox
boundary carried as explicit \emph{false} capability fields in every
receipt, so that no consumer can honestly read a receipt as
hermetic-execution evidence.
\item \emph{Exact-byte-backed claim and theorem registries} whose
loaders reject every structural irregularity and whose records carry
deterministic digests bound to the registry bytes, so that parsed
content cannot be substituted while raw identity is retained.
\end{enumerate}

\subsection{Unrepresentable versus fail-closed}\label{ch01:unrep}

The inversion has two distinct enforcement modes, and keeping them
separate sharpens both the design and its review. A behavior is
\emph{unrepresentable} when no value of any kernel type expresses it:
there is no way to write down an identity link between references in
different namespaces, because the construction that would express it
rejects at creation, and no other construction exists. Unrepresentability
is the stronger mode---the dangerous state is not guarded against but
absent---and it is available only where the kernel owns the vocabulary,
which is why the semantic-identity layer leans on it hardest.

A behavior is \emph{fail-closed} when it is representable as an attempt
but the attempt is rejected at a boundary, with an error, before any
downstream artifact exists. Canonicalization is the model case: a caller
can certainly pass a float, a set, or a cyclic structure to the encoder;
the encoder's job is to refuse loudly rather than guess. Fail-closed is
the right mode wherever the kernel receives input from outside its own
type system---bytes from disk, values from callers, records from
ledgers---because there the dangerous state arrives from a vocabulary
the kernel does not control.

The review consequence is worth stating: unrepresentability is reviewed
by reading type definitions and constructors (is there really no path?),
while fail-closure is reviewed by reading rejection tests (does every
bad input class have a fixture that is refused?). The bounded receipts
cited later in the monograph speak to the second mode directly---the
suite's rejection cases are enumerated---and to the first mode only via
constructor tests. Both reviews remain open obligations in the
independent sense; the monograph reports what the receipts exercised,
which is not the same thing.

\subsection{The constant-false authorization stance}\label{ch01:constfalse}

Threaded through all five mechanisms is the single commitment already
named: structural consistency is never authorization. In the v0.1
candidate this commitment takes its most extreme available form. There is
\emph{no code path that grants authorization}. The runtime decision
model's field \texttt{authorization\_granted}, and the compatibility
spelling \texttt{allowed}, are constant false; claim-bearing intent is
denied even when a locally constructed ledger contains the text
\textsc{licensed}, because the kernel has no authenticated reviewer and
no legal trust root, and a self-hash or a Python-local sentinel is not
issuance provenance. The most a nonclaim research operation can obtain is
a structural replay comparison---a statement that two hash-verified
executions agreed---which is deliberately typed as evidence, not as
permission.

It is worth being precise about what this stance does and does not
assert, because ``the kernel always says no'' can be misread as either
paranoia or theater. The stance does not assert that authorization is
impossible in principle, that review can never license a claim, or that
the program intends to remain forever unlicensed. It asserts that
\emph{this} artifact, lacking a trust root, is not the place where
authorization can originate, and that encoding this fact as a constant is
strictly safer than encoding it as a configurable check. A configurable
check can be misconfigured, mocked in testing and left mocked, or
satisfied by laundered inputs; a constant with no writing code path can
only be changed by changing the kernel, which is a reviewed act. The
governing sources also state minimum design constraints that any future
claim-authorization implementation must bind---an exact trusted
registry/schema manifest, an authenticated authority and trust root,
reviewer and admission receipts, evidence-packet identities and
verification status, exact claim identity, and the code, model,
environment, semantic, and calibration identities of the requested
operation---and until a separately reviewed design binds all of them,
the kernel's answer remains no.

\subsection{What the inversion costs}\label{ch01:costs}

Honesty about the thesis requires stating its price. A fail-closed kernel
is inconvenient in exactly the ways it is safe. Rejection-first
canonicalization means ordinary Python values---floats, sets, ad hoc
dicts with non-string keys---do not silently flow into evidence; every
producer must adapt to the encoding, not the reverse. Disjoint
namespaces mean cross-ontology analysis requires constructing explicit
link objects with sources and notes, which is more work than typing the
same string in two columns. Constant-false authorization means the kernel
cannot, in its current form, be placed on any route that needs a yes.
The design accepts these costs deliberately, on the argument---developed
across the threat analyses of the later chapters---that each convenience
foregone corresponds to a historically observed failure mode, and that
the program has already paid, in voided work products, the larger price
of the convenient alternative. Whether the balance is struck correctly
is a question for the independent review that the kernel's own status
block says remains open.

\subsection{Why a kernel, and why now}\label{ch01:whykernel}

One further design question deserves an answer in the introduction,
because it shapes how every later chapter should be read: why is this a
\emph{kernel}---a small, separate artifact---rather than a discipline
applied inside the program's existing codebases?

The first reason is reviewability. The properties the kernel claims are
of the kind that can only be established by reading everything: an
unrepresentability claim is a claim about the absence of a path, and
absence is audited by exhaustion. Exhaustive reading is feasible for a
handful of small modules with no dependencies on the semantics they
police; it is not feasible for a research monorepo. The kernel's
smallness is therefore not modesty but a precondition of its central
claims. Every capability added to the kernel dilutes this property,
which is one more reason capabilities arrive only through reviewed
gates.

The second reason is dependency direction. The kernel is designed as
the shared substrate of a set of companion research candidates---board
modeling, typed physical quantities, a finite typed-arrow kernel, and a
categorical evidence layer among them---and the dependency arrow is
strict: companions consume the kernel's contracts; the kernel depends
on none of their semantics. This directionality is what lets the kernel
be reviewed in isolation and what prevents the classic failure in which
the integrity layer quietly imports the domain layer and every audit of
the former becomes an audit of both. A discipline diffused through the
codebases would have no such direction; it would be everywhere and
auditable nowhere.

The third reason is timing. The program is at a stage where a future
CAS v8 is a candidate, not a codebase---which means the frame can still
be built \emph{before} the semantics that will inhabit it. Retrofitting
identity and evidence discipline onto a living system means competing
with every existing convenience and migrating every existing artifact;
building the frame first means the first v8 artifact born is born
inside it. The two historical findings of this chapter are, on this
reading, the program's tuition payments: the conflation incident and
the self-hashed receipts each demonstrated, at the cost of voided work,
what the absence of the frame permits. Building the kernel now, while
there is nothing yet to overclaim about, is the cheap moment.

None of this argues that the kernel is sufficient---the open-obligation
set says plainly that it is not---only that it is correctly shaped and
correctly timed for the narrow job it claims: to be the part of the
system that everything else must go through, small enough to read,
strict enough to trust as far as structural properties go, and honest
enough to say ``no'' about everything beyond them.

\section{Scope and non-goals}\label{ch01:scope}

The kernel's credibility depends on the exactness of its boundary at
least as much as on its mechanisms, so we state the boundary here in
full and will restate the relevant portions wherever a chapter's content
might otherwise invite an inflated reading.

\begin{boundary}[Not a rewrite of CAS v7.7]\label{ch01:bd-notrewrite}
CAS Core 8 is deliberately not a rewrite of CAS v7.7. The v7.7 framework
---its operators, domains, drives, registers, notation, and canon---
continues to exist and evolve on its own track. The kernel neither
reimplements v7.7 semantics nor supersedes any v7.7 artifact. Its only
contact with legacy semantics is the compatibility profile, which is
explicitly \texttt{legacy\_only}: it executes hash-verified legacy
compiler bytes for comparison purposes and produces no CAS v8 semantic
translation. No v8 parser and no full semantic migration is claimed
anywhere in the candidate.
\end{boundary}

\begin{boundary}[Not on production routes]\label{ch01:bd-noproduction}
The kernel is not installed on any production route. Nothing a user
touches is served by, gated through, or dependent on CAS Core 8 in its
v0.1 form. The kernel is a research candidate whose correct venue is
review, and its receipts describe test and research executions only.
\end{boundary}

\begin{boundary}[No hermetic execution, no semantic correctness, no authority]\label{ch01:bd-noclaims}
Three specific overclaims are disclaimed globally, here and throughout.
First, the compatibility profile achieves Python-process isolation, not
an operating-system sandbox; the host's native libraries and operating
system are not content-addressed, and every receipt carries these
limitations as explicit false capability fields. No statement in this
monograph is a claim of hermetic execution. Second, nothing in the
kernel establishes the semantic correctness of any CAS content---the
kernel checks identity, encoding, and structural consistency, never
meaning. Third, nothing in the kernel or in this monograph constitutes
claim authority: admission state is \textsc{open}, product state is
\textsc{not licensed}, disclosure is withheld, and independent software,
formal-methods, semantics, and product-safety review remains an open
obligation.
\end{boundary}

Beyond these three, the governing sources enumerate further open
obligations that scope the kernel from the other side: a hermetic
subprocess and toolchain bundle is not established by Python-level
isolation alone; Python/Rust engine conformance is not established and
the engine profiles remain separate pending a differential-testing
campaign; no import-DAG manifest for multi-file programs is admitted,
so multi-file legacy programs are outside the compatibility profile; and
no claim-authority implementation is admitted. Where a later chapter
touches one of these, it will cite the obligation rather than quietly
narrowing it.

\begin{remark}[Relation to the wider program]\label{ch01:relation}
Because the kernel polices vocabulary drawn from the wider program's
canon, it is easy to misread it as adjudicating that canon. It does
not. The authoritative definitions of the program's semantic
families---which texts define which elements, which names are
canonical, which correspondences are positional rather than
identical---are fixed by the program's canon documents and their own
governance, entirely outside the kernel. The kernel consumes the
canon's facts as inputs: a reference's definition digest is the digest
of bytes the canon supplies; a crosswalk's non-identity typing encodes
a distinction the canon draws. If the canon changes, references bound
to superseded definitions become visibly stale, which is the intended
behavior---but the kernel has no opinion about what the canon should
say, and nothing in this monograph should be read as a gloss on, or a
rival authority to, any canon document. The division of labor is
exact: the canon owns meaning; the kernel owns the bookkeeping that
keeps meaning from being silently swapped.
\end{remark}

Two scope remarks about the monograph itself. First, the monograph's
evidential basis is fixed: the candidate specification, the v0.2
research paper that supersedes the v0.1 white paper, the kernel
implementation sources, and the bounded Gate-4 receipts they cite. The
monograph adds exposition, organization, and argument; it adds no facts.
Where the exposition needs a fact the sources do not supply, the text
says the sources do not supply it. Second, the monograph inherits the
kernel's claim discipline. Words like ``verified,'' ``guaranteed,'' and
``proved'' are used only for properties enforced by construction or
exercised by the cited receipts, always with their fixture-bound
character stated; design intentions are called design intentions.

\section{Objections and replies}\label{ch01:objections}

Three objections recur when the kernel's stance is presented, and
answering them here sharpens what the monograph is and is not
arguing.

\emph{Objection: the kernel is bureaucracy---the real work is the
semantics, and this is paperwork around an empty center.} The reply
is the program's own ledger. The failures that voided work were not
semantic errors that better semantics would have caught; they were
frame errors---conflated identities, promoted receipts---that
corrupted sound work. Paperwork is the pejorative for provenance
maintained after the fact; the kernel is the opposite bet, that
identity and evidence discipline maintained \emph{by construction}
is what lets future semantic work accumulate instead of dissolving
under audit. Whether the bet pays is an empirical question the
program will answer; that the alternative was already tried and
failed is a matter of record.

\emph{Objection: constant-false authorization is theater---if the
kernel can never say yes, nothing real depends on it, so its
refusals cost nothing and prove nothing.} Half of this is correct
and accepted: in v0.1, no production behavior depends on the gate,
which is precisely the claim of
Boundary~\ref{ch01:bd-noproduction}. What the objection misses is
where the value of a constant-false period lies: it is the interval
in which every \emph{other} part of the kernel---identities,
encodings, receipts, registries---hardens under real use while the
authorization question is held safely trivial. The alternative
sequencing, standing up a configurable yes/no gate first and
hardening its inputs later, is exactly the shape of system in which
the historical receipts incident occurred. The constant is not the
end state; the stated future-authorization constraints are. The
constant is the refusal to improvise the end state.

\emph{Objection: the fail-closed posture will not survive contact
with users---friction accumulates, and someone will build the
convenient bypass.} This is the serious one, and the monograph does
not pretend the kernel is immune; a bypass built beside the kernel
is outside its boundary and no type system prevents it. Two things
can be said. Structurally, the kernel minimizes the bypass surface
by being a substrate rather than a checkpoint: companions consume
its contracts as their data model, so bypassing it means forgoing
the data model, not just skipping a check. Institutionally, the
kernel makes the bypass \emph{visible}: artifacts that do not carry
its typed identities and boundaries are identifiable as
non-conforming at review. Beyond that, the objection names a
governance problem, and the kernel's honest position is that
governance problems are not solvable by kernels---which is, after
all, its thesis.

\section{Reading guide}\label{ch01:guide}

The monograph is organized so that a reader can enter at the level of
generality they need and descend.

\emph{The background chapter} that follows is a tutorial survey of the
standard material the kernel builds on: cryptographic hashing and
content addressing; canonical serialization and the reasons it is
genuinely hard, with the JSON Canonicalization Scheme as the reference
point; reproducible builds and supply-chain integrity; provenance
models built on entities, activities, and derivations; the
integrity-versus-authorization distinction as security engineering has
long drawn it; threat-modeling method; and the distinction between
process isolation and sandboxing. None of that material is project
novelty, and the chapter is explicit about this: the kernel's relation
to the background is adaptation and composition, not invention. Readers
fluent in these topics can skim the chapter for its comparison tables
and its closing map from background concepts to kernel mechanisms.

\emph{The semantic-identity chapter} develops the first mechanism in
full: the five-field structural reference and the rationale for each
field, with particular attention to definition-hash binding as the
detector for definitional drift; the disjoint namespace catalog; the
link-kind system for cross-ontology correspondences and the historical
incident that motivates its refusal of identity links; the
injective-rendering invariant with its design argument and its
delimiter fixtures; the exact-type equality rule and the spoofing
scenario it blocks; and worked examples of legal and illegal reference
usage, grounded in the reference implementation.

\emph{Subsequent chapters} treat, in order: the research-boundary
object and its recheck discipline; the canonical evidence encoding as a
closed rule set with its threat analysis; the legacy compiler
compatibility profile, its deterministic-closure protocol, its
concurrency and failure disciplines, and its non-sandbox boundary; the
byte-backed claim and theorem registries; the runtime decision model
and the constant-false authorization invariant, together with the
stated design constraints on any future authorization; the consolidated
threat analysis pairing each historically observed failure with the
mechanism that now blocks it; the kernel's relation to its companion
candidates and the strict direction of dependency that keeps the kernel
reviewable in isolation; the bounded software evidence, stated with its
exact limits, including the fact that the repository-wide collector was
not clean at the bound receipt and that the focused suites therefore
support no repository-wide claim; and the complete open-obligation
set, which functions as the monograph's honest ending: a list of the
things this kernel does not yet establish, maintained with the same
care as the list of things it does.

Three reading paths may help three kinds of reader. A reviewer assessing
the kernel's safety story should read this introduction, the
semantic-identity chapter, and the threat-analysis material, then audit
the boundary statements of every other chapter. An engineer intending
to build on the kernel's contracts should read the background chapter's
canonicalization sections and then the mechanism chapters in order,
treating every fail-closed rejection rule as part of the interface. A
reader from the wider research program who mainly needs to know what
the kernel does \emph{not} do can read
Sections~\ref{ch01:promotion}--\ref{ch01:scope} of this chapter and the
open-obligation material, and take away the one sentence the whole
design is built to keep true: nothing in this kernel, and nothing in
this monograph, converts a receipt into an authority.

\subsection{How to read the evidence statements}\label{ch01:evidencehow}

Because later chapters cite the bounded Gate-4 receipts repeatedly, the
reader should calibrate once, here, what those citations can and cannot
carry. The receipts record, at a fixed date and commit, a focused suite
of sixty-one passing test methods over the kernel, identified by a
digest of the suite's test identifiers, within an integrated run of two
hundred eight methods and two hundred seventy-two parameterized
subtests over fifty-four hash-bound artifacts. Three limits govern
every citation of these numbers. First, they are \emph{historical}:
they describe runs that happened, not a standing property of the
codebase; any change to the code or fixtures after the receipt is
outside their reach. Second, they are \emph{fixture-bound}: a rejection
test establishes that the tested bad input was refused, and supports an
inductive expectation about its class only to the degree the class is
well represented---a judgment the receipts themselves cannot make.
Third, they are \emph{scope-bound}: at the bound receipt the
repository-wide test collector was not clean (four pre-existing
package-path namespace collisions after some twelve hundred collected
items), so the focused suites support no repository-wide verification
claim, and the monograph will never use them to make one. When a later
chapter writes ``exercised by the receipts,'' it means exactly this
much and no more. The reader who internalizes this calibration once
can then read the mechanism chapters without the text having to
re-hedge every sentence---though the boundary environments will
re-state the limits wherever an inflated reading would be most
tempting, on the theory that redundancy about limits is the one form
of redundancy this subject rewards.

It is also worth saying what would count as \emph{more} evidence,
since the monograph's claim discipline can otherwise read as a counsel
of despair. Independent re-execution of the suites from the hash-bound
artifacts would upgrade the receipts from reported to reproduced.
Independent review of the type definitions and constructors would
address the unrepresentability claims, which no test can fully carry.
Adversarial fixture construction---an external attempt to produce
colliding renderings, type-punned encodings, or laundered
authorizations---would test the fail-closed surfaces where they matter.
And the separately stated review obligations (software,
formal-methods, semantics, product-safety) name the institutional
version of all three. The kernel's stance is not that such evidence is
unobtainable; it is that until obtained, the kernel and its monograph
say so.

\begin{remark}[Terminology]\label{ch01:terminology}
Throughout the monograph, ``the kernel'' means the CAS Core 8 v0.1
research candidate as fixed by its specification and implementation
sources; ``the program'' means the surrounding TKS/CAS research effort;
``receipt'' means a structural, machine-generated record of observed
software behavior; ``canon'' means the program's authoritative
definitional documents; and ``Ring 4'' abbreviates the joint boundary
condition---admission \textsc{open}, product \textsc{not licensed},
public disclosure withheld---that every artifact discussed here
carries. Where a term of art from the background literature is used, it
is defined at first use in the background chapter.
\end{remark}
